\documentclass[10pt,a4paper]{amsart}
\usepackage[margin=19mm]{geometry}
\usepackage{amsmath,amssymb,mathtools}
\usepackage[scr=rsfs,cal=euler]{mathalfa}
\usepackage{xfrac}
\usepackage[unicode,hidelinks,bookmarks=false]{hyperref}
\newcommand{\ie}{i.e.\ }
\newcommand{\cf}{cf.\ }
\newcommand{\resp}{respectively\ }
\newcommand{\Subsection}{\subsection}
\providecommand{\nobreakdash}{\nobreak\hbox{-}}
\setlength{\emergencystretch}{2em}
\multlinegap0pt
\usepackage{amssymb}
\usepackage[shortlabels]{enumitem}
\usepackage{tikz}
\newtheorem{theorem}{Theorem}
\newtheorem{remark}{Remark}
\newtheorem{lemma}{Lemma}
\newtheorem{corollary}{Corollary}
\newcommand{\id}{\operatorname{id}}
\title{The Borel Distinguishing Number of Schreier Graphs}
\author{Junhao Chen, Jie Zou}
\date{Source: arXiv:2609.15003v1 (CC BY 4.0)}
\begin{document}
\maketitle

\noindent\emph{Source and licence.} The Borel Distinguishing Number of Schreier Graphs. Version arXiv:2609.15003v1 (CC BY 4.0), distributed by arXiv under the Creative Commons Attribution 4.0 International licence, http://creativecommons.org/licenses/by/4.0/, as recorded in the arXiv licence metadata for this exact version and shown on the abstract page https://arxiv.org/abs/2609.15003v1. Full licence notice: \texttt{licenses/arxiv-2609.15003-CC-BY-4.0.txt}.\par\smallskip

\noindent\textit{Editorial scope.} Counted unit: the article's theorem upperbound (source0.tex lines 359--361). The article's complete original proof of it is delivered with it (source0.tex lines 362--422, one \texttt{proof} environment of 61 lines). No mathematics is written, repaired or re-derived: the statement and the proof are the article's own lines, in the article's own notation, and no retained line is substituted or re-flowed.  Retained uncounted as the source-local context the proof needs: lines 133--136; lines 202--204; lines 332--335; lines 350--357.  Nothing was omitted from any retained span, so the package carries no omission record.  No retained line contains a citation, so the wrapper carries no bibliography and no bibliography entry.  No retained line contains a figure environment, an image-inclusion command or drawing code, so no image is delivered and the package declares no asset dependency.  Editorial formatting only, in addition: an AMS \texttt{amsart} preamble that restates the article's own declarations and macros used by the retained lines, namely the environments \texttt{theorem}, \texttt{remark}, \texttt{lemma}, \texttt{corollary} on separate counters, together with the standard packages those lines need.  The retained environments print in this document as Theorem 1, Remark 1, Lemma 1, Corollary 1 in order of appearance, whereas the article prints the same items with the numbering of its own full document.  The complete native source of the article as distributed in the arXiv e-print for this exact version is bundled unchanged as \texttt{sources/210383/source0.tex}.
\begin{theorem}
\label{KST}
For any Borel graph $\mathcal{G}$ with bounded maximum degree $\Delta(\mathcal{G})$, its Borel chromatic number satisfies $\chi_B(\mathcal{G})\le \Delta(\mathcal{G}) + 1$.    
\end{theorem}

\begin{remark}\label{11}
By the proposition above, if $f$ is a Borel distinguishing $Y$-coloring, then the symbolic factor map $\varphi_f$ is a Borel equivariant embedding into $Y^\Gamma$. Also, to find a Borel distinguishing $Y$-coloring $f$, it suffices to find a Borel equivariant injection $\varphi: X \longrightarrow Y^\Gamma$ with the following property: for any $\gamma \in A_e(\Gamma) \setminus \{\id\}$ and $x,y\in X$, $\varphi(x) \ne \gamma\cdot\varphi(y)$, since $f(x)=\varphi(x)(e)$ will be the desired Borel distinguishing $Y$-coloring. 
\end{remark}

\begin{lemma}\label{Ae}
For every 
$\gamma\in A_e(\mathbb{Z}^d)$, there exist a permutation $\sigma$ of $\{1,\dots,d\}$ and a vector $\varepsilon=(\varepsilon_1,\ldots,\varepsilon_d)\in\{-1,1\}^d$ such that for each $x\in \mathbb{Z}^d$, $\gamma(x) $ can be written as $\sum_{i=1}^d \varepsilon_i x_i e_{\sigma(i)}.$
\end{lemma} 

\begin{corollary}\label{cube}
Let $\gamma\in A_e(\mathbb{Z}^d)$, and let
$B_g=Q_k+g$ be a cube of even length $k$ centered at
$g\in\mathbb{Z}^d$. Then $\gamma(B_g)=Q_k+\gamma(g).$ Precisely, $\gamma(\partial B_g)=\partial\gamma(B_g),
\gamma(M(B_g))=M(\gamma(B_g)),\gamma(C(B_g))=C(\gamma(B_g)).$
Furthermore, for every $u\in Q_k$,
$\gamma(u+g)=\gamma(u)+\gamma(g).$
\end{corollary}

\begin{theorem}\label{upperbound}
Let $d\ge 1, n\ge 2$ be two integers, and let $\mathcal{G}=(X,E)$ be the Schreier graph induced by the free part of the shift action $\mathbb{Z}^d \curvearrowright n^{\mathbb{Z}^d}$. Then $D_{B}(\mathcal{G})\le n+1$.
\end{theorem}

\begin{proof}
Our proof relies on  Remark  \ref{11}. We will pointwise construct a Borel equivariant injection $\varphi: X\longrightarrow (n+1)^{\mathbb{Z}^d}$ such that $\varphi(\alpha)\ne\gamma\cdot\varphi(\beta)$ for all $\alpha,\beta\in X$ and all $\gamma\in A_e(\mathbb{Z}^d)\setminus\{\id\}$. 

We first choose a sufficiently large even number $k$ such that:
\begin{enumerate}
    \item $k>$ Max$\{6,3d+2\}$.
    \item $n^{(k+1)^d}\leq (n+1)^{(k-3)^d}$.
\end{enumerate}
We  show that there is a pattern $p$ such that, for
every  $\gamma\in A_e(\mathbb{Z}^d)\setminus\{id\}$, every
$\alpha\in (n+1)^{\mathbb{Z}^d}$, and every cube $B$ of length
$k$, if $\alpha$ has pattern $p$ on $B$, then
$\gamma\cdot\alpha$ does not have pattern $p$ on
$\gamma(B)$.
 
If $d\ge 2$, let
$v=(1,2,\dots,d-1,\frac{k-2}{2})\in M(Q_k)$. If $d=1$, let
$v=\frac{k-2}{2}$. We will show that $\gamma(v)\neq v$ for every 
$\gamma\in A_e(\mathbb{Z}^d)\setminus\{id\}$. Suppose $\gamma(v)=v$.
By Lemma \ref{Ae}, there exist $\sigma$ and
$\varepsilon\in\{-1,1\}^d$ such that $\gamma(v)=\sum_{i=1}^d \varepsilon_i v_i e_{\sigma(i)}.$ Hence $\varepsilon_i v_i=v_{\sigma(i)}$ for every $i$. Since the coordinates of $v$ are positive and pairwise distinct, we have $\varepsilon_i=1$ and $\sigma(i)=i$
for every $i$. Thus $\gamma=\id$, a contradiction. We color $v$ with $1$ and
all other points in $M(Q_k)$ with $0$. This is the pattern
$p$ we need.

Then we show that, for each cube $B$ of length $k$, if a
coloring has pattern $p$ on $B$, then, after applying any
 $\gamma\in A_e(\mathbb{Z}^d)\setminus\{id\}$, the induced pattern
on $\gamma(B)$ is not equal to $p$. Let $B=Q_k+g$. The unique
point colored with $1$ in $M(B)$ is $v+g$. By
Corollary \ref{cube},  the unique point colored with $1$ in  $M(\gamma(B))$ is $\gamma(v)+\gamma(g)$. Since
$\gamma(v)\neq v$ for every 
$\gamma\in A_e(\mathbb{Z}^d)\setminus\{id\}$, this  pattern is
different from $p$.



By our choice of $k$, there exists an injection $\psi: n^{Q_k} \longrightarrow (n+1)^{C(Q_k)}$ . We have also fixed a pattern $p$ of length $k$ for this $k$.

By Theorem \ref{KST}, there is a Borel proper coloring $c:X\longrightarrow l$ of the power graph $\mathcal{G}^{4dk}$, where $l\in\mathbb{N}$ is an integer. For each $\alpha\in X$, define the center set
\[
\mathcal{C}(\alpha)=\left\{g\in\mathbb{Z}^d:c(g\cdot\alpha)=\min_{h\in\mathbb{Z}^d}c(h\cdot\alpha)\right\}.
\]
We define $\varphi$ as follows: for each $g\in \mathcal{C}(\alpha)$, let $B_g$ denote the cube of length $k$ centered at $g$. We recolor this cube according to the following rules:
\begin{enumerate}
    \item Fill $\partial B_g$ with the new color $n$ to help us identify those changed cubes.
    \item On $M(B_g)$, we place the pattern $p$ to break the symmetry.
    \item For $C(B_g)$, we shift $B_g$ to the standard cube $Q_k$, encode the original coloring of $B_g$ into its core using $\psi$,
and then shift it back to its original position.
\end{enumerate}

We do this for all $g\in \mathcal{C}(\alpha)$ and leave all other areas with their original colors. Let the result be $\varphi(\alpha)$. 

Now we show that $\varphi$ satisfies our requirements. We only need to check these changed areas. Because $c$ and the action of ${\mathbb{Z}}^d$ are Borel and, for a fixed $x\in \mathbb{Z}^d$, determining whether $x$ lies in a changed area involves a countable verification, $\varphi$ is Borel. For the equivariance, note that for each $g\in\mathbb{Z}^d$, 
\[\forall x\in {\mathbb{Z}}^d, x\in \mathcal{C}(g\cdot\alpha)\iff c((x+g)\cdot\alpha)\text{ has the least color}\iff x+g\in \mathcal{C}(\alpha).\]
 Hence $\mathcal{C}(g\cdot\alpha)=\mathcal{C}(\alpha)-g$, which means all the centers move equivariantly, as do all the cubes. Moreover, the color pattern does not change under the shift action. This implies equivariance.

 For injectivity, since $c$ is a proper coloring, any two distinct centers are at a distance greater than $4dk$. So any two distinct changed cubes $B$, $B'$ are at a distance greater than $dk$. Then we can identify each changed cube through its boundary of length $k$. In fact, the color $n$ does not occur in the middle layers or outside these cubes. Thus, even if the color $n$ occurs inside a core, the core is not large enough to have length $k$, so we will not confuse the core with the boundary. After we identify the position of each changed cube,  we can check the core and restore the original coloring through the injection $\psi$.
 
Finally, we show that for all $\alpha,\beta\in X$ and all $\gamma\in A_e(\mathbb{Z}^d)\setminus\{\mathrm{id}\}$, $\varphi(\alpha)\ne\gamma\cdot\varphi(\beta)$. Suppose there are $\alpha,\beta\in X$ and $\gamma\in A_e(\mathbb{Z}^d)\setminus\{\mathrm{id}\}$ such that $\varphi(\alpha)=\gamma\cdot\varphi(\beta)$. We fix a changed cube $B$ in $\varphi(\beta)$. By our choice of $p$, $\gamma\cdot\varphi(\beta)$ has another pattern $p'$ on the image cube. Note that $\gamma(B)$ is a changed cube in $\varphi(\alpha)$. However, by rule (2), the pattern $p'$ will not occur on the middle layer of any changed cube, a contradiction.
\end{proof}

\end{document}
